% Figure 1: the headline. (a) a schematic of what .s1 does differently; (b-g) every value read from generated/*.tsv and
% generated/*_boxes.tex (written by `python -m s1bench report`). One tikzpicture with the axes on a grid: b and c share
% their rows (one per format), d is square, e and f share their rows (one per loader) under b and c, g sits under d.
% \figbox{pos}{q1}{median}{q3}{low whisker}{high whisker}{colour}: one horizontal box, in axis coordinates.
\newcommand{\figbox}[7]{%
  \draw[#7!80!black, line width=0.5pt] (axis cs:#5,#1) -- (axis cs:#2,#1);
  \draw[#7!80!black, line width=0.5pt] (axis cs:#4,#1) -- (axis cs:#6,#1);
  \draw[#7!80!black, fill=#7, fill opacity=0.18, line width=0.5pt] (axis cs:#2,{#1-0.32}) rectangle (axis cs:#4,{#1+0.32});
  \draw[#7!70!black, line width=1.2pt] (axis cs:#3,{#1-0.32}) -- (axis cs:#3,{#1+0.32});}
\begin{figure*}[tbp]
\centering
\pgfmathsetlengthmacro{\figT}{max(\linewidth,18.1cm)}   % design width: panel a needs 18.1 cm; scaled to fit below
\pgfmathsetlengthmacro{\figY}{2.15cm}             % y tick labels of b and e
\pgfmathsetlengthmacro{\figGap}{0.9cm}            % between b and c, e and f (shared rows)
\pgfmathsetlengthmacro{\figD}{4.9cm}              % side of the square panel d (as placed by report_r3.FIG1D_SIDE)
\pgfmathsetlengthmacro{\figDG}{1.5cm}             % before d and g: their y labels
\pgfmathsetlengthmacro{\figW}{(\figT-\figY-\figGap-\figDG-\figD-0.25cm)/2}   % width of b, c, e and f: rows fill \figT
\ifdefined\IEEEtransversionmajor                  % TCBB: the figure and caption fill a page
  \pgfmathsetlengthmacro{\figH}{4.2cm}            % height of e, f and g
  \pgfmathsetlengthmacro{\figV}{2.0cm}            % between rows: x labels below, titles above
  \pgfmathsetlengthmacro{\figA}{5.6cm}            % top of the figure to the top of b and c
\else
  \pgfmathsetlengthmacro{\figH}{3.0cm}
  \pgfmathsetlengthmacro{\figV}{1.75cm}
  \pgfmathsetlengthmacro{\figA}{5.2cm}
\fi
\begin{adjustbox}{max width=\linewidth}%
\begin{tikzpicture}[
  sbox/.style={draw=gen!80!black, fill=gen!12, rounded corners=1.5pt, inner sep=3.5pt, align=center,
    font=\sffamily\scriptsize, text=ink2},
  obox/.style={draw=sone!85!black, fill=sone!12, rounded corners=1.5pt, inner sep=3.5pt, align=center,
    font=\sffamily\scriptsize, text=ink2},
  note/.style={font=\sffamily\tiny, text=ink2, align=center},
  arr/.style={-{Stealth[length=3.5pt]}, thick, ink2!80},
  rowlbl/.style={font=\sffamily\small\bfseries, anchor=west},
]
\pgfplotsset{every axis title/.append style={at={(0,1)}, anchor=south west, xshift=-4pt, text depth=0.3ex, text height=1.6ex}}
% ======================= (a) what .s1 does differently =======================
% Left: the same six counts of one cell, stored the standard way (gene index and count as 32-bit numbers, then a byte
% compressor) and the .s1 way (genes numbered most-detected first, steps between them, one symbol stream, a model).
% Right: reading a minibatch, one chunk per round trip against the index first and every block at once.
\node[anchor=north west, font=\sffamily\small\bfseries, text=black, inner sep=0pt] at (0,0) {a\quad What \sone{} does differently};
\node[font=\sffamily\scriptsize\bfseries, text=black] at (5.6,-0.6) {1. Describe the counts, not just their bytes};
\node[font=\sffamily\scriptsize\bfseries, text=black, align=center] at (15.0,-0.6) {2. Fetch a minibatch's cells in one round of requests};
\draw[line width=0.4pt, gen!40] (11.65,-0.4) -- (11.65,-4.15);
\node[rowlbl, text=gen!80!black] at (0,-1.55) {CSC};
\node[rowlbl, text=sone!85!black] at (0,-3.15) {\sone};
\tikzset{
  cel/.style={draw=#1!85!black, fill=#1!12, minimum width=0.27cm, minimum height=0.34cm, inner sep=0pt,
    font=\sffamily\tiny, text=ink2},
  wcel/.style={draw=gen!80!black, fill=gen!18, minimum width=0.48cm, minimum height=0.34cm, inner sep=0pt,
    font=\sffamily\tiny, text=ink2},
  alb/.style={font=\sffamily\tiny, text=gen!90!black, anchor=east, inner sep=1pt},
  stg/.style={font=\sffamily\tiny, text=ink2, align=center, anchor=north},
  hl/.style={anchor=west, font=\sffamily\scriptsize, text=ink2, align=left},
}
\begin{scope}[xshift=-0.5cm]
% --- storing, CSC: each count as a 32-bit gene index and a 32-bit count, then a byte compressor ---
\node[alb] at (2.25,-1.33) {gene};
\node[alb] at (2.25,-1.77) {count};
\foreach \g [count=\i from 0] in {41,103,220,778,912,1554} \node[wcel] at ({2.54+0.48*\i},-1.33) {\g};
\foreach \v [count=\i from 0] in {1,1,3,2,1,1} \node[wcel] at ({2.54+0.48*\i},-1.77) {\v};
\node[stg] at (3.74,-2.0) {two 32-bit numbers per count};
\draw[arr] (5.3,-1.55) -- (8.95,-1.55);
\node[sbox, anchor=west, minimum width=1.55cm] (gz) at (9.0,-1.55) {gzip, zstd\\[-1pt]{\tiny compress bytes}};
\draw[arr] (gz.east) -- (10.85,-1.55);
\node[hl] at (10.9,-1.55) {\textbf{\nm{size-h5ad_gzip}}{\tiny $^\ast$}\\[-1pt]bits/count};
% --- storing, .s1: genes numbered by how often they are detected, steps between them, one symbol stream, a model ---
\node[alb] at (2.25,-2.93) {gene no.};
\node[alb] at (2.25,-3.37) {count};
\foreach \g [count=\i from 0] in {0,2,3,5,9,10} \node[cel=sone] at ({2.44+0.27*\i},-2.93) {\g};
\foreach \v [count=\i from 0] in {1,1,3,2,1,1} \node[cel=dec] at ({2.44+0.27*\i},-3.37) {\v};
\node[stg] at (3.11,-3.6) {genes numbered,\\most detected first};
\draw[arr] (4.0,-3.15) -- (4.35,-3.15);
\foreach \s [count=\i from 0] in {0,1,0,1,3,0} \node[cel=sone] at ({4.58+0.27*\i},-2.93) {\s};
\foreach \v [count=\i from 0] in {1,1,3,2,1,1} \node[cel=dec] at ({4.58+0.27*\i},-3.37) {\v};
\node[stg] at (5.25,-3.6) {step: genes skipped\\since the last one};
\draw[arr] (6.14,-3.15) -- (6.49,-3.15);
\foreach \s/\c [count=\i from 0] in {0/sone,1/sone,0/sone,3/dec,1/sone,2/dec,3/sone,0/sone} \node[cel=\c] at ({6.72+0.27*\i},-3.15) {\s};
\node[stg] at (7.67,-3.6) {one stream: each step, then\\the count if it is above 1};
\draw[arr] (8.82,-3.15) -- (8.97,-3.15);
\node[obox, anchor=west, minimum width=1.55cm] (md) at (9.0,-3.15) {short codes\\[-1pt]{\tiny from a model of}\\[-2pt]{\tiny the counts}};
\draw[arr] (md.east) -- (10.85,-3.15);
\node[hl] at (10.9,-3.15) {\textbf{\nm{size-s1}}\\[-1pt]bits/count};
\end{scope}
% --- reading, CSC: one round trip after another (fetch arrow, then decode box) ---
\def\pX{11.95}\def\rtA{0.72}\def\dcA{0.36}\def\cyc{1.16}
\foreach \k in {0,1,2} {
  \pgfmathsetmacro{\xs}{\pX+\k*\cyc}
  \draw[gen!80!black, line width=0.5pt, -{Stealth[length=3pt]}] (\xs,-1.55) -- ++(\rtA,0);
  \fill[gen!55, rounded corners=0.8pt] ({\xs+\rtA+0.02},-1.39) rectangle ++(\dcA,-0.32);
}
\node[note, anchor=west, align=left] at ({\pX+3*\cyc+0.3},-1.55) {one chunk at a time: fetch,\\wait a round trip, decode,\\repeat; or download the\\whole file first};
% --- reading, .s1: the index first, then every block of a minibatch at once, decoded on every core ---
\fill[sone!85!black, rounded corners=0.8pt] (\pX,-2.82) rectangle ++(0.3,-0.66);
\node[note, anchor=south] at ({\pX+0.15},-2.78) {index};
\foreach \l in {0,1,2,3} {
  \draw[sone!85!black, line width=0.5pt, -{Stealth[length=3pt]}] ({\pX+0.4},{-2.9-0.165*\l}) -- ++(\rtA,0);
  \fill[sone!60, rounded corners=0.6pt] ({\pX+0.42+\rtA},{-2.835-0.165*\l}) rectangle ++(0.28,-0.13);
}
\node[note, anchor=west, align=left] at ({\pX+3*\cyc+0.3},-3.15) {read the index once, then\\fetch all of a minibatch's\\blocks at once and decode\\them on every core};
\draw[gen!60, densely dotted, line width=0.4pt] ({\pX+0.42+\rtA+0.33},-2.7) -- ({\pX+0.42+\rtA+0.33},-3.62);
\draw[gen!60, densely dotted, line width=0.4pt] ({\pX+3*\cyc},-1.15) -- ({\pX+3*\cyc},-3.62);
\draw[-{Stealth[length=3pt]}, gen!70, line width=0.4pt] (\pX,-3.75) -- ({\pX+3*\cyc+0.05},-3.75) node[below left, note, text=gen, inner sep=1pt] {time};
\node[note, anchor=north west, text=gen!90!black] at (0,-4.05) {$^\ast$\,gzip-compressed AnnData, median over \nm{nMatrices} matrices};
% ======================= (b) size: compression ratio of every matrix =======================
\begin{axis}[name=pb, at={(\figY,-\figA)}, anchor=north west, scale only axis, width=\figW, height=\figD,
  title={b\quad Smallest files}, title style={xshift={1pt-\figY}},
  xmin=0, xmax=17, xtick={0,4,...,16}, xlabel={compression against 32-bit CSC ($\times$)},
  ytick=data, ytick style={draw=none}, yticklabels from table={\gen{fig1a_size.tsv}}{label}, y dir=reverse,
  ymin=-0.6, ymax=15.6, axis x line*=bottom, axis y line*=left, clip=false]
\addplot[draw=none, forget plot] table[x=ratio_median, y=pos] {\gen{fig1a_size.tsv}};   % sets the y ticks
\draw[gen!60, line width=0.4pt, dashed] (axis cs:1,-0.5) -- (axis cs:1,15.5);
\input{\gen{fig1a_boxes.tex}}
\expandafter\pgfplotsinvokeforeach\expandafter{\otherformats,s1}{
  \addplot[only marks, mark=*, mark size=0.75pt, c-#1, opacity=0.6, discard if not={format}{#1}] table[x=ratio, y=y] {\gen{fig1a_points.tsv}};
}
\end{axis}
% ======================= (c) speed: read rate of every matrix on four cores =======================
\begin{axis}[name=pc, at={($(pb.north east)+(\figGap,0)$)}, anchor=north west, scale only axis, width=\figW, height=\figD,
  title={c\quad Fastest to read},
  xmin=0, xmax=260, xtick={0,50,...,250},
  xlabel={million non-zeros/s, 4 cores},
  ytick={0,...,15}, yticklabels={}, ytick style={draw=none}, y dir=reverse, ymin=-0.6, ymax=15.6,
  axis x line*=bottom, axis y line*=left, clip=false]
\input{\gen{fig1_speed_boxes.tex}}
\expandafter\pgfplotsinvokeforeach\expandafter{\otherformats,s1}{
  \addplot[only marks, mark=*, mark size=0.75pt, c-#1, opacity=0.6, discard if not={format}{#1}] table[x=mnnzs, y=y] {\gen{fig1_speed_points.tsv}};
}
\end{axis}
% ======================= (d) size against speed (figures/fig1b_panel.tex) =======================
\def\fbLog{0}\def\fbPts{1}\def\fbTitle{d\quad Smaller and faster at once}\def\fbInPicture{1}
\def\fbAxisOpts{name=pd, at={($(pc.north east)+(\figDG,0)$)}, anchor=north west, scale only axis, width=\figD, height=\figD}
% Fig. 1b: stored size against cold read rate, from generated/fig1c_tradeoff.tsv (medians and quartiles per format) and
% generated/fig1b_points.tsv (every matrix). Set before \input: \def\fbLog{1} for log axes or {0} for linear ones, and
% \def\fbPts{1} to draw every matrix (with \def\fbPts{0} each median gets its interquartile ranges instead), and the
% panel title in \fbTitle.
\ifnum\fbLog=1
  \def\fbopts{xmode=log, ymode=log, xmin=3, xmax=110, ymin=2, ymax=2000,
    xtick={5,10,20,50,100}, xticklabels={5,10,20,50,100}, ytick={10,100,1000}, yticklabels={10,100,1000},
    xlabel={stored bits per non-zero (log scale)}, ylabel={cold read rate, million non-zeros/s (log scale)}}
  \def\fbiso{25,50,100,200}\def\fbxlo{3}\def\fbtop{1500}\def\fbanchor{anchor}
\else
  \def\fbopts{xmin=0, xmax=70, ymin=0, ymax=260, xtick={0,10,...,70}, ytick={0,50,...,250},
    xlabel={stored bits per non-zero}, ylabel={cold read rate, million non-zeros/s}}
  \def\fbiso{25,50,100}\def\fbxlo{0.2}\def\fbtop{250}\def\fbanchor{anchor_lin}
\fi
% Inside another tikzpicture (Fig. 1), \fbAxisOpts gives the axis position and size and no picture is opened here.
\def\fbaxis#1{\begin{axis}[#1]}
\ifdefined\fbAxisOpts\else\begin{tikzpicture}\def\fbAxisOpts{width=0.47\linewidth, height=7.6cm}\fi
\edef\fbAll{\unexpanded\expandafter{\fbopts},\unexpanded\expandafter{\fbAxisOpts}}
\expandafter\fbaxis\expandafter{\fbAll,
  title={\fbTitle}, clip mode=individual,
  axis x line*=bottom, axis y line*=left,
  lbl/.style={nodes near coords, point meta=explicit symbolic, visualization depends on={value \thisrow{\fbanchor}\as\lblanchor},
    every node near coord/.append style={font=\sffamily\tiny, anchor=\lblanchor, text=ink2, inner sep=1.5pt}},
  iqr/.style={error bars/x dir=both, error bars/y dir=both, error bars/x explicit, error bars/y explicit,
    error bars/error bar style={line width=0.5pt, opacity=0.8}, error bars/error mark=none},
}
% iso-lines of Eq. (1) at B = 1 Gbit/s: R_eff = R  <=>  D = 1/(1/R - b/1000), D and R in million non-zeros/s, b in bits;
% each ends, labeled, at the top of the plot (the unit only on the 25 line, so the close labels on the left fit)
\expandafter\pgfplotsinvokeforeach\expandafter{\fbiso}{
  \addplot[domain=\fbxlo:{1000*(1/#1-1/\fbtop)}, samples=150, dashed, sone!55, forget plot] {1/(1/#1 - x/1000)};
  \node[anchor=south west, font=\sffamily\tiny, text=sone!80!black, inner sep=0.5pt]
    at (axis cs:{1000*(1/#1-1/\fbtop)},\fbtop) {#1\ifnum#1=25\,M/s\fi};
}
\node[anchor=north east, font=\sffamily\tiny, text=sone!80!black, align=right, inner sep=1pt] at (rel axis cs:0.99,0.89)
  {dashed: non-zeros per second\\delivered over a 1\,Gbit/s link};
% labels: beside each median on log axes; on linear axes beside, above or below each point, placed by
% generated/fig1d_labels.tex to avoid overlaps
\tikzset{fblabel/.style={font=\sffamily\tiny, text=ink2, inner sep=0.6pt, outer sep=0pt, fill=white, fill opacity=0.8, text opacity=1}}
\ifnum\fbLog=1\pgfplotsset{fblbl/.style={lbl}}\else\pgfplotsset{fblbl/.style={}}\input{\gen{fig1d_labels.tex}}\fi
\ifnum\fbPts=1
\expandafter\pgfplotsinvokeforeach\expandafter{\otherformats,s1}{
  \addplot[only marks, mark=*, mark size=0.7pt, c-#1, opacity=0.4, forget plot, discard if not={format}{#1}] table[x=bits, y=mnnzs] {\gen{fig1b_points.tsv}};
}
\pgfplotsset{fbmed/.style={draw=white, line width=0.4pt}}
\else
\pgfplotsset{fbmed/.style={iqr}}
\fi
\expandafter\pgfplotsinvokeforeach\expandafter{\otherformats}{
\addplot[only marks, fblbl, mark=*, mark size=1.8pt, c-#1, fbmed, discard if not={format}{#1}]
  table[x=bits, y=mnnzs, meta=label, x error plus expr=\thisrow{bits_q3}-\thisrow{bits}, x error minus expr=\thisrow{bits}-\thisrow{bits_q1},
    y error plus expr=\thisrow{mnnzs_q3}-\thisrow{mnnzs}, y error minus expr=\thisrow{mnnzs}-\thisrow{mnnzs_q1}] {\gen{fig1c_tradeoff.tsv}};
}
\addplot[only marks, fblbl, mark=*, mark size=2.4pt, c-s1, fbmed, discard if not={format}{s1}]
  table[x=bits, y=mnnzs, meta=label, x error plus expr=\thisrow{bits_q3}-\thisrow{bits}, x error minus expr=\thisrow{bits}-\thisrow{bits_q1},
    y error plus expr=\thisrow{mnnzs_q3}-\thisrow{mnnzs}, y error minus expr=\thisrow{mnnzs}-\thisrow{mnnzs_q1}] {\gen{fig1c_tradeoff.tsv}};
\end{axis}
\ifdefined\fbInPicture\else\end{tikzpicture}\let\fbAxisOpts\undefined\fi%

% ======================= (e, f) a PyTorch training loader at 1 Gbit/s =======================
\pgfplotsset{fig1mb/.style={xbar, bar shift=0pt, bar width=0.5, ymin=-0.5, ymax=3.5, y dir=reverse,
  scale only axis, width=\figW, height=\figH, axis x line*=bottom, axis y line*=left, clip=false,
  ytick={0,1,2,3}, ytick style={draw=none},
  error bars/error bar style={ink2, line width=0.6pt}, error bars/error mark options={rotate=90, mark size=2pt, ink2}},
  fig1val/.style={nodes near coords, point meta=explicit, every node near coord/.append style={font=\sffamily\tiny,
  anchor=west, text=ink2, inner sep=2pt}}}
\begin{axis}[xmode=log, log origin=infty, fig1mb, name=pe, at={($(pb.south west)+(0,-\figV)$)}, anchor=north west,
  title={e\quad Loads first PyTorch minibatch fastest}, title style={xshift={1pt-\figY}},
  xmin=0.02, xmax=40, xtick={0.1,1,10}, xticklabels={0.1,1,10}, xlabel={seconds (log scale)},
  yticklabels from table={\gen{fig1_mb.tsv}}{label}]
\pgfplotsinvokeforeach{s1,zarr3,zarr,soma}{
  \addplot[fill=c-#1, draw=none, discard if not y={fmt}{#1}, error bars/.cd, x dir=both, x explicit]
    table[x=first, y=pos, x error plus expr=\thisrow{first_hi}-\thisrow{first}, x error minus expr=\thisrow{first}-\thisrow{first_lo}] {\gen{fig1_mb.tsv}};
}
\addplot[draw=none, fig1val, nodes near coords={\pgfmathprintnumber[fixed, fixed zerofill, precision=2]{\pgfplotspointmeta}\,s}]
  table[x=first_hi, y=pos, meta=first] {\gen{fig1_mb.tsv}};
\end{axis}
\begin{axis}[xmode=log, log origin=infty, fig1mb, name=pf, at={($(pe.north east)+(\figGap,0)$)}, anchor=north west,
  title={f\quad and later minibatches},
  xmin=2, xmax=4000, xtick={10,100,1000}, xticklabels={10,100,1000}, xlabel={milliseconds (log scale)}, yticklabels={}]
\pgfplotsinvokeforeach{s1,zarr3,zarr,soma}{
  \addplot[fill=c-#1, draw=none, discard if not y={fmt}{#1}, error bars/.cd, x dir=both, x explicit]
    table[x=batch_ms, y=pos, x error plus expr=\thisrow{batch_hi}-\thisrow{batch_ms}, x error minus expr=\thisrow{batch_ms}-\thisrow{batch_lo}] {\gen{fig1_mb.tsv}};
}
\addplot[draw=none, fig1val, nodes near coords={\pgfmathprintnumber[fixed, precision=0]{\pgfplotspointmeta}\,ms}]
  table[x=batch_hi, y=pos, meta=batch_ms] {\gen{fig1_mb.tsv}};
\end{axis}
% ======================= (g) read rate against cores =======================
\begin{axis}[name=pg, at={($(pf.north east)+(\figDG,0)$)}, anchor=north west, scale only axis, width=\figD, height=\figH,
  title={g\quad Best parallel efficiency},
  xmin=0, xmax=46, ymin=0, ymax=14, xtick={1,8,16,24,32,44}, ytick={0,2,...,14},
  xlabel={cores}, ylabel={speed-up over one core},
  axis x line*=bottom, axis y line*=left,
  legend style={font=\sffamily\tiny, draw=none, fill=white, fill opacity=0.85, text opacity=1, at={(1.0,1.0)},
    anchor=north east, legend columns=2, cells={anchor=west}, row sep=-2pt, column sep=3pt},
  legend image post style={xscale=0.55}]
\addplot[gen!70, dashed, line width=0.6pt, domain=1:14, samples=2, forget plot] {x};
\node[anchor=north west, font=\sffamily\tiny, text=gen!90!black, inner sep=1.5pt] at (axis cs:8,8) {1:1};
\addplot[name path=sq3, draw=none, forget plot, discard if not={format}{s1}] table[x=threads, y=sp_q3] {\gen{fig1_scaling.tsv}};
\addplot[name path=sq1, draw=none, forget plot, discard if not={format}{s1}] table[x=threads, y=sp_q1] {\gen{fig1_scaling.tsv}};
\addplot[c-s1, fill opacity=0.18, draw=none, forget plot] fill between[of=sq3 and sq1];
\addplot[c-s1, very thick, mark=*, mark size=1.5pt, discard if not={format}{s1}] table[x=threads, y=speedup] {\gen{fig1_scaling.tsv}};
\addlegendentry{\sone}
\addplot[c-soma, thick, mark=*, mark size=1.0pt, discard if not={format}{soma}] table[x=threads, y=speedup] {\gen{fig1_scaling.tsv}};
\addlegendentry{TileDB-SOMA}
\addplot[c-zarr3, thick, mark=*, mark size=1.0pt, discard if not={format}{zarr3}] table[x=threads, y=speedup] {\gen{fig1_scaling.tsv}};
\addlegendentry{Zarr v3}
\addplot[c-parquet_coo, thick, mark=*, mark size=1.0pt, discard if not={format}{parquet_coo}] table[x=threads, y=speedup] {\gen{fig1_scaling.tsv}};
\addlegendentry{Parquet}
\addplot[c-mtx_gz, thick, mark=*, mark size=1.0pt, discard if not={format}{mtx_gz}] table[x=threads, y=speedup] {\gen{fig1_scaling.tsv}};
\addlegendentry{MTX.gz}
\addplot[c-csc_i32, thick, mark=*, mark size=1.0pt, discard if not={format}{csc_i32}] table[x=threads, y=speedup] {\gen{fig1_scaling.tsv}};
\addlegendentry{CSC, int32}
\addplot[c-h5ad_gzip, thick, mark=*, mark size=1.0pt, discard if not={format}{h5ad_gzip}] table[x=threads, y=speedup] {\gen{fig1_scaling.tsv}};
\addlegendentry{h5ad gzip}
\end{axis}
\end{tikzpicture}%
\end{adjustbox}
\caption{\textbf{\sone{} files are the smallest and, on four cores, the fastest to read.}
\textbf{(a)} The two ideas behind \sone. Left: six counts of one cell, stored as compressed sparse column (CSC) arrays
with a byte compressor, as in gzip-compressed AnnData, and as \sone{} does, with genes numbered from most to least often detected, the steps between them and the counts merged into one
stream of small symbols, and short codes from a model of the counts. Right: reading a minibatch. Fig.~\ref{fig:how} shows
the layout.
\textbf{(b)} Compression of each format's file (counts, gene identifiers and barcodes) against uncompressed CSC arrays with
32-bit indices and values, one point per matrix (\nm{nMatrices} matrices); boxes, quartiles and median; whiskers, the most
extreme matrices within 1.5 interquartile ranges.
\textbf{(c)} Read rate of the same files from a cold page cache into an in-memory sparse matrix on \nm{threadsHeadline}
cores, rows as in (b).
\textbf{(d)} Medians of (b) and (c) as bits per non-zero against read rate; faint points, single matrices. Dashed lines
join formats that deliver the same number of non-zeros per second over a 1\,Gbit/s link (Eq.~\ref{eq:reff}); up and to
the left is better.
\textbf{(e, f)} PyTorch training loaders reading 512-cell minibatches of \nm{fig1b-dataset} from an emulated object store
at 1\,Gbit/s, each decoded into a sparse tensor: (e) seconds from a new process to the first minibatch, median and range
of three seeds; (f) milliseconds for each later minibatch, median and 5th--95th percentiles of \nm{mb3-batches}
minibatches from three seeds.
\textbf{(g)} Speed-up of the read rate from a cold file over one core against cores on a 48-core node, median over
\nm{scale-s1-n} matrices of each matrix's own speed-up; band, interquartile range of \sone; dashed gray, the 1:1 line of
perfect scaling. Shown are the formats whose readers can use several cores, with uncompressed CSC and gzip-compressed
AnnData for reference.}
\label{fig:wow}
\end{figure*}
